% BHU PYQ -- Transcribed & Corrected
% Paper: HINMD-23: Adivasi Jeevan Evam Sahitya (आदिवासी जीवन एवं साहित्य)
% Exam: Under Graduate Semester II Examination, 2024-25 (NEP)
% Category: Multidisciplinary Course (MDC)
% Department: Hindi / Faculty of Arts

\documentclass[11pt,a4paper]{article}
\usepackage[a4paper,top=2.5cm,bottom=2.5cm,left=2.5cm,right=2.5cm,headheight=14pt]{geometry}
\usepackage{amsmath,amssymb}
\usepackage[T1]{fontenc}
\usepackage[utf8]{inputenc}
\usepackage{lmodern,microtype,fancyhdr,enumitem,hyperref,lastpage,parskip}

\hypersetup{colorlinks=true,urlcolor=black,linkcolor=black,pdftitle={HINMD23: Adivasi Jeevan Evam Sahitya -- BHU Sem II 2024-25 (NEP MDC)}}
\pagestyle{fancy}\fancyhf{}
\fancyhead[L]{\small \textit{Banaras Hindu University}}
\fancyhead[R]{\small \textit{HINMD-23: Adivasi Jeevan Evam Sahitya (2024-25)}}
\fancyfoot[C]{\small Page \thepage\ of \pageref{LastPage}}

\newcommand{\pts}[1]{\hfill \textit{[#1]}}

\begin{document}

\begin{center}
  {\large\bfseries BANARAS HINDU UNIVERSITY}\\[2pt]
  {\normalsize Under Graduate Semester II Examination, 2024-25 (NEP)}\\[6pt]
  \rule{0.8\linewidth}{0.5pt}\\[6pt]
  {\large\bfseries FACULTY OF ARTS / DEPARTMENT OF HINDI}\\[4pt]
  {\large\bfseries Multidisciplinary Course (MDC)}\\[4pt]
  {\large\bfseries Paper Code: HINMD-23 --- आदिवासी जीवन एवं साहित्य (Adivasi Jeevan Evam Sahitya)}\\[4pt]
  {\normalsize Time: 2:30 Hours\hfill Maximum Marks: 70}
\end{center}

\rule{\linewidth}{0.4pt}\smallskip

\noindent \textit{अनुदेश : निर्देशानुसार प्रश्नों के उत्तर दीजिए। प्रश्नों के अंक उनके सम्मुख अंकित हैं।}

\bigskip

\noindent \textbf{1.} निम्नलिखित दीर्घउत्तरीय प्रश्नों में से किन्हीं तीन प्रश्नों के उत्तर दीजिए। \pts{3 $\times$ 10 = 30}

\begin{enumerate}[label=(\roman*)]
  \item आदिवासी साहित्य की अवधारणा को समझाइए।
  \item आदिवासी साहित्य की प्रमुख प्रवृत्तियों को स्पष्ट कीजिए।
  \item आदिवासी साहित्य और प्रकृति के परस्पर संबंध को रेखांकित कीजिए।
  \item 'वनकन्या' कहानी का उद्देश्य स्पष्ट कीजिए।
  \item 'भंवर' कहानी के प्रमुख पात्रों का चरित्र चित्रण कीजिए।
  \item पठित उपन्यास के आधार पर आदिवासी समाज की चुनौतियों को बताइए।
  \item पठित कविताओं के आधार पर आदिवासी जीवन संघर्ष व उसकी विडम्बनाओं को समझाइए।
  \item पठित कहानियों के आधार पर आदिवासी साहित्य में व्यक्त आदिवासी जीवन दर्शन व परम्पराओं को स्पष्ट कीजिए।
\end{enumerate}

\bigskip

\noindent \textbf{2.} निम्नलिखित काव्यांशों में से किन्हीं चार की सप्रसंग व्याख्या कीजिए। \pts{4 $\times$ 6 = 24}

\begin{enumerate}[label=(\roman*)]
  \item तप्त हृदय की न्यारी आह\\
  बन के अतितम घोर प्रदाह।\\
  झुलसने चली देखिए चाह,\\
  दाह के मन-मंदिर की थाह।।

  \item सुगम बना उपहार हार है ,\\
  किस के उर पर झूलेगा,\\
  किस प्रेमी की प्रणय चाह को,\\
  झूम-झूम कर चूमेगा?

  \item बच्चों के लिए मैदान\\
  पशुओं के लिए हरि-हरि घास\\
  बूढ़ों के लिए पहाड़ों की शांति,\\
  और इस अविश्वास भरे दौर में\\
  थोड़ा सा विश्वास\\
  थोड़ी सी उम्मीद\\
  थोड़े से सपने

  \item और उसके हाथ में मत देना मेरा हाथ\\
  जिसके हाथों ने कभी कोई पेड़ नहीं लगाए\\
  फसलें नहीं उगाईं जिन हाथों ने\\
  जिन हाथों ने दिया नहीं कभी किसी का साथ\\
  किस का बोझ नहीं उठाया

  \item अभी अभी\\
  सुन्न हुई उसकी देह से\\
  बिजली की लपलपाती कौंध निकली\\
  जेल की दीवार लांघती\\
  तीर की तरह जंगलों में पहुंची\\
  एक-एक दरख्त, बेल, झुरमुट\\
  पहाड़, नदी, झरना\\
  वन प्राणी-पखेरू, कीट, सरीसृप\\
  खेत- खलिहान, बस्ती\\
  वहाँ की हवा, धूल, जमीन में समा गई।

  \item मैं जंगल का पुश्तैनी दावेदार हूँ\\
  पुरखे और उनके दावे मरते नहीं\\
  मैं भी मर नहीं सकता\\
  मुझे कोई भी\\
  जंगलों से बेदखल नहीं कर सकता\\
  उलगुलान ! उलगुलान !! उलगुलान !!!
\end{enumerate}

\bigskip

\noindent \textbf{3.} निम्नलिखित में से किन्हीं चार पर टिप्पणी कीजिए : \pts{4 $\times$ 4 = 16}

\begin{enumerate}[label=(\roman*)]
  \item आदिवासी साहित्य की परिभाषा
  \item खरगोशों का कष्ट कहानी का सार
  \item बडी दीदी कहानी का मूल भाव
  \item रचने होंगे ग्रंथ कविता की मूल संवेदना
  \item उतनी दूर बिहाना मत बाबा कविता का मूल स्वर
  \item मिनाम उपन्यास की पात्र यामी का चरित्र चित्रण
\end{enumerate}

\bigskip
\begin{center}
  \textbf{*****}
\end{center}

\end{document}
